\chapter{Neural tangent kernel and linearization near initialization}

\section{Introduction}

This chapter formalizes the theory of networks near their random initialization,
corresponding to Chapter~4 of the deep learning theory notes (Telgarsky 2021).
The central idea is to compare a neural network $f(x;W)$ to its
\emph{first-order Taylor approximation} $f_0(x;W)$ at a random initialization $W_0$.
The approximation $f_0$ is affine in $W$ while remaining nonlinear in $x$,
which greatly eases analysis.

\begin{enumerate}
  \item Section~\ref{sec:ntk:setup} fixes the network architecture (a scaled shallow network)
        and defines the Taylor linearization at initialization.
  \item Section~\ref{sec:ntk:linearization} shows that $f - f_0$ shrinks to zero as the
        width $m \to \infty$, for both smooth activations (via Taylor's theorem) and for the
        ReLU (via a Gaussian concentration argument).
  \item Section~\ref{sec:ntk:kernel} defines the \emph{neural tangent kernel} (NTK) as the
        Gram matrix of the gradient features at initialization, derives its limit as
        $m \to \infty$ by the strong law of large numbers, and gives a closed form for the ReLU.
  \item Section~\ref{sec:ntk:universal} shows that the associated RKHS is a universal
        approximator over a suitably chosen domain.
\end{enumerate}

Throughout, we write $\|\cdot\|_{\mathrm{F}}$ for the Frobenius norm of a matrix,
$\langle A, B \rangle_{\mathrm{F}} = \mathrm{tr}(A^\top B)$ for the Frobenius inner product,
and $\|x\|$ for the Euclidean norm of a vector.
All probability statements are with respect to the standard Gaussian initialization
(Definition~\ref{def:ntk:gaussianInit}).


\section{Shallow network setup}\label{sec:ntk:setup}

We study a one-hidden-layer (``shallow'') network with input dimension $d \in \mathbb{N}$,
width $m \in \mathbb{N}$, fixed outer layer $a \in \mathbb{R}^m$, and variable inner weight
matrix $W \in \mathbb{R}^{m \times d}$.

\begin{definition}[Scaled shallow network]\label{def:ntk:shallowNetwork}
  \lean{NTK.ShallowNetwork}
Let $\sigma : \mathbb{R} \to \mathbb{R}$ be an activation function, $d, m \in \mathbb{N}$,
and fix $a \in \mathbb{R}^m$ with $|a_j| \le 1$ for all $j \in [m]$.
The \emph{scaled shallow network} is the function
\[
  f : \mathbb{R}^d \times \mathbb{R}^{m \times d} \to \mathbb{R}, \qquad
  f(x; W) := \frac{1}{\sqrt{m}} \sum_{j=1}^{m} a_j \, \sigma(w_j^\top x),
\]
where $w_j \in \mathbb{R}^d$ denotes the $j$-th row of $W$.
\end{definition}

\begin{remark}
The $1/\sqrt{m}$ normalization ensures that the empirical NTK defined in
Section~\ref{sec:ntk:kernel} is an average over neurons and has a finite limit as $m \to \infty$.
We hold $a$ fixed and only optimize $W$; this is the simplest setup that gives a
non-convex but tractable training problem.
\end{remark}

\begin{definition}[Gaussian random initialization]\label{def:ntk:gaussianInit}
  \lean{NTK.gaussianInit}
The \emph{standard Gaussian initialization} is the product probability measure on
$\mathbb{R}^{m \times d}$ under which the rows $w_1, \ldots, w_m \in \mathbb{R}^d$ are
drawn i.i.d.\ from $\mathcal{N}(0, I_d)$.
We write $W_0 \sim \mathcal{N}(0, I_d)^{\otimes m}$ for a draw from this initialization.
\end{definition}

\begin{definition}[First-order Taylor linearization]\label{def:ntk:linearization}
  \uses{def:ntk:shallowNetwork}
  \lean{NTK.linearization}
Let $W_0 \in \mathbb{R}^{m \times d}$ be a fixed initialization.
The \emph{first-order Taylor linearization} of $f(\cdot\,; \cdot)$ at $W_0$ is
\[
  f_0(x; W) :=
    f(x; W_0) + \langle \nabla_W f(x; W_0),\, W - W_0 \rangle_{\mathrm{F}},
\]
where the gradient $\nabla_W f(x; W_0) \in \mathbb{R}^{m \times d}$ has rows
$a_j \sigma'(w_{0,j}^\top x)\, x^\top / \sqrt{m}$.

For an activation $\sigma$ that is differentiable almost everywhere (e.g.\ the ReLU),
this expands to
\[
  f_0(x; W) = \frac{1}{\sqrt{m}} \sum_{j=1}^{m} a_j \Bigl[
    \sigma(w_{0,j}^\top x) - \sigma'(w_{0,j}^\top x)\,w_{0,j}^\top x
    + \sigma'(w_{0,j}^\top x)\,w_j^\top x\Bigr].
\]
\end{definition}

\begin{remark}[ReLU simplification]\label{rmk:ntk:reluLinearization}
  \uses{def:ntk:linearization}
For the ReLU $\sigma_{\mathrm{ReLU}}(z) = \max\{0,z\}$, the identity
$\sigma_{\mathrm{ReLU}}(z) = z\,\sigma_{\mathrm{ReLU}}'(z)$ (holding for a.e.\ $z$)
eliminates the constant term, giving the cleaner form
\[
  f_0(x; W) = \frac{1}{\sqrt{m}} \sum_{j=1}^{m} a_j\, \sigma'(w_{0,j}^\top x)\, w_j^\top x
  = \langle \nabla_W f(x; W_0),\, W \rangle_{\mathrm{F}}.
\]
\end{remark}


\section{Networks near initialization are almost linear}\label{sec:ntk:linearization}

\subsection{Smooth activations}

For smooth activations, the linearization error follows immediately from Taylor's theorem
and requires no probabilistic argument.

\begin{definition}[$\beta$-smooth activation]\label{def:ntk:betaSmooth}
  \lean{NTK.BetaSmooth}
An activation $\sigma : \mathbb{R} \to \mathbb{R}$ is \emph{$\beta$-smooth} if $\sigma$ is
twice differentiable everywhere and $|\sigma''(z)| \le \beta$ for all $z \in \mathbb{R}$.
\end{definition}

\begin{proposition}[Smooth linearization bound]\label{prop:ntk:smoothLinearization}
  \uses{def:ntk:shallowNetwork, def:ntk:linearization, def:ntk:betaSmooth}
  \lean{NTK.smoothLinearizationBound}
Suppose $\sigma$ is $\beta$-smooth and $|a_j| \le 1$ for all $j \in [m]$.
Let $W_0 \in \mathbb{R}^{m \times d}$ be a fixed initialization.
For any $x \in \mathbb{R}^d$ with $\|x\| \le 1$ and any $W \in \mathbb{R}^{m \times d}$,
\[
  |f(x; W) - f_0(x; W)| \le \frac{\beta}{2\sqrt{m}} \|W - W_0\|_{\mathrm{F}}^2.
\]
\end{proposition}

\begin{proof}
  \uses{def:ntk:betaSmooth, def:ntk:shallowNetwork, def:ntk:linearization}
By Taylor's theorem applied to $\sigma$ with remainder,
\[
  |\sigma(r) - \sigma(s) - \sigma'(s)(r-s)| \le \frac{\beta(r-s)^2}{2}.
\]
Summing over neurons, using $|a_j| \le 1$ and $|\langle w_j - w_{0,j}, x\rangle| \le \|w_j - w_{0,j}\| \cdot \|x\| \le \|w_j - w_{0,j}\|$:
\begin{align*}
  |f(x;W) - f_0(x;W)|
  &= \frac{1}{\sqrt{m}}\left|\sum_{j=1}^m a_j
      \bigl[\sigma(w_j^\top x) - \sigma(w_{0,j}^\top x) - \sigma'(w_{0,j}^\top x)(w_j-w_{0,j})^\top x\bigr]
     \right| \\
  &\le \frac{1}{\sqrt{m}} \sum_{j=1}^m \frac{\beta(w_j^\top x - w_{0,j}^\top x)^2}{2} \\
  &\le \frac{\beta}{2\sqrt{m}} \sum_{j=1}^m \|w_j - w_{0,j}\|^2
   = \frac{\beta}{2\sqrt{m}} \|W - W_0\|_{\mathrm{F}}^2. \qedhere
\end{align*}
\end{proof}

\begin{remark}
Proposition~\ref{prop:ntk:smoothLinearization} does not require Gaussian initialization:
it holds for any $W, W_0$.
The key consequence is that as $m \to \infty$, the linearization error $|f - f_0|$ decays
as $\mathcal{O}(1/\sqrt{m})$ whenever $\|W - W_0\|_{\mathrm{F}} = \mathcal{O}(1)$.
\end{remark}

\subsection{ReLU linearization via Gaussian concentration}

The ReLU is not smooth, so we instead exploit the Gaussian initialization $W_0$:
with high probability, $|w_{0,j}^\top x|$ is large, so a perturbation of bounded
Frobenius norm does not change the sign of $w_j^\top x$ for most $j$.

To make this precise, we first bound the probability that a single neuron's pre-activation falls into an ambiguous region where its sign might flip.

\begin{lemma}[Gaussian pre-activation probability bound]\label{lem:ntk:probSignAmbiguous}
  \lean{NTK.prob_signAmbiguous_le_tau}
For $x \in \mathbb{R}^d$ with $\|x\| > 0$ and $\tau > 0$,
\[
  \Pr_{w \sim \mathcal{N}(0, I_d)}\!\bigl[ |w^\top x| \le \tau\|x\| \bigr] \le \tau.
\]
\end{lemma}

We then use a standard Hoeffding bound for sums of independent indicator variables:

\begin{lemma}[Hoeffding bound for indicators]\label{lem:ntk:hoeffding}
  \lean{NTK.hoeffding_indicators_pi}
Let $P_1, \ldots, P_m$ be independent indicator variables with $\mathbb{E}[P_j] \le p$.
With probability at least $1-\delta$, $\sum_{j=1}^m P_j \le mp + \sqrt{\frac{m}{2}\ln\frac{1}{\delta}}$.
\end{lemma}

Combining these gives the sign concentration lemma for the entire network:

\begin{lemma}[Sign concentration under Gaussian initialization]\label{lem:ntk:signConcentration}
  \uses{def:ntk:gaussianInit, lem:ntk:probSignAmbiguous, lem:ntk:hoeffding}
  \lean{NTK.reluSignConcentration}
Let $x \in \mathbb{R}^d$ with $\|x\| > 0$ and let $W_0 \sim \mathcal{N}(0, I_d)^{\otimes m}$.
For any $\tau > 0$ and $\delta \in (0, 1)$, with probability at least $1 - \delta$ over $W_0$,
\[
  \#\!\bigl\{j \in [m] : |w_{0,j}^\top x| \le \tau\|x\|\bigr\}
  \;\le\; m\tau + \sqrt{\tfrac{m}{2} \ln \tfrac{1}{\delta}}.
\]
\end{lemma}

\begin{proof}
  \uses{def:ntk:gaussianInit}
For each $j \in [m]$, define the indicator $P_j := \mathbf{1}[|w_{0,j}^\top x| \le \tau\|x\|]$.
By rotational invariance of $\mathcal{N}(0,I_d)$, we may assume $\|x\| = 1$ and
$w_{0,j}^\top x \sim \mathcal{N}(0,1)$, so
\[
  \mathbb{E}[P_j] = \Pr[|Z| \le \tau] \le \frac{2\tau}{\sqrt{2\pi}} \le \tau, \quad Z \sim \mathcal{N}(0,1).
\]
The variables $(P_j)_{j=1}^m$ are i.i.d.\ Bernoulli, so Hoeffding's inequality gives
\[
  \Pr\!\left[\sum_j P_j > m\tau + \sqrt{\tfrac{m}{2}\ln\tfrac{1}{\delta}}\right] \le \delta. \qedhere
\]
\end{proof}

We also need to bound the number of neurons that move significantly from their initialization.

\begin{lemma}[Large perturbation bound]\label{lem:ntk:cardLargePerturb}
  \lean{NTK.card_largePerturb_bound}
Let $W_0, W \in \mathbb{R}^{m \times d}$ and $r > 0$. The number of neurons $j$ for which $\|w_j - w_{0,j}\| \ge r$ is bounded by $\|W - W_0\|_{\mathrm{F}}^2 / r^2$.
\end{lemma}

Let $S$ be the union of the set of neurons with ambiguous signs at initialization (bounded by Lemma~\ref{lem:ntk:signConcentration}) and those that move significantly (bounded by Lemma~\ref{lem:ntk:cardLargePerturb}). Outside of $S$, the signs are preserved:

\begin{lemma}[Sign preservation outside bad set]\label{lem:ntk:signPreserved}
  \lean{NTK.sign_preserved_outside_badSet, NTK.relu_error_eq_zero_outside_badSet}
If $|w_{0,j}^\top x| > r\|x\|$ and $\|w_j - w_{0,j}\| < r$, then $\operatorname{sgn}(w_j^\top x) = \operatorname{sgn}(w_{0,j}^\top x)$. Consequently, the Taylor expansion of the ReLU activation is exact, and the linearization error for this neuron is exactly zero.
\end{lemma}

Thus, the total linearization error sums only over the "bad" set $S$, which we can bound using Cauchy-Schwarz:

\begin{lemma}[Cauchy-Schwarz bound over bad set]\label{lem:ntk:csBound}
  \lean{NTK.relu_error_cs_bound}
For any index set $S \subseteq [m]$, the error is bounded by:
\[
  \frac{1}{\sqrt{m}} \sum_{j \in S} \left| \sigma(w_j^\top x) - \sigma(w_{0,j}^\top x) - \sigma'(w_{0,j}^\top x)(w_j - w_{0,j})^\top x \right| \le \frac{\|W - W_0\|_{\mathrm{F}}}{\sqrt{m}} \sqrt{|S|}.
\]
\end{lemma}

These ingredients combine into the main ReLU linearization bound.

\begin{lemma}[ReLU linearization bound]\label{lem:ntk:reluLinearization}
  \uses{def:ntk:shallowNetwork, def:ntk:linearization, def:ntk:gaussianInit,
        lem:ntk:signConcentration, lem:ntk:cardLargePerturb, lem:ntk:signPreserved, lem:ntk:csBound}
  \lean{NTK.reluLinearizationBound}
  \lean{NTK.reluLinearizationBound_secondOrder}
Let $\sigma = \sigma_{\mathrm{ReLU}}$, let $W_0 \sim \mathcal{N}(0,I_d)^{\otimes m}$,
and fix $B \ge 0$ and $x \in \mathbb{R}^d$ with $\|x\| \le 1$.
With probability at least $1 - \delta$ over $W_0$, uniformly over all
$W \in \mathbb{R}^{m \times d}$ with $\|W - W_0\|_{\mathrm{F}} \le B$,
\[
  |f(x; W) - f_0(x; W)| \le \frac{2B^{4/3} + B\ln(1/\delta)^{1/4}}{m^{1/6}}.
\]
Moreover, for every $V \in \mathbb{R}^{m \times d}$ with $\|V - W_0\|_{\mathrm{F}} \le B$,
\[
  |f(x; V) - (f(x; W) + \langle \nabla_W f(x; W),\, V - W\rangle_{\mathrm{F}})| \le
  \frac{6B^{4/3} + 3B\ln(1/\delta)^{1/4}}{m^{1/6}}.
\]
\end{lemma}

\begin{proof}
  \uses{lem:ntk:signConcentration, def:ntk:gaussianInit, def:ntk:shallowNetwork}
Set $r := B^{2/3}/m^{1/3}$ and define the ``bad'' index sets
\begin{align*}
  S_1 &:= \bigl\{j \in [m] : |w_{0,j}^\top x| \le r\|x\|\bigr\}, \\
  S_2 &:= \bigl\{j \in [m] : \|w_j - w_{0,j}\| \ge r\bigr\}, \\
  S   &:= S_1 \cup S_2.
\end{align*}
By Lemma~\ref{lem:ntk:signConcentration} with $\tau = r$,
\[
  |S_1| \le rm + \sqrt{\tfrac{m}{2}\ln\tfrac{1}{\delta}} \quad \text{w.p.\ } \ge 1 - \delta.
\]
By the Frobenius norm bound on $W - W_0$, $|S_2| \le B^2/r^2$.
Choosing $r = B^{2/3}/m^{1/3}$ balances the two terms:
\[
  |S| \le m^{2/3}\!\left(2B^{2/3} + \sqrt{\ln(1/\delta)}\right).
\]
For $j \notin S$, both $|w_{0,j}^\top x| > r\|x\|$ and $\|w_j - w_{0,j}\| < r$, so
$\operatorname{sgn}(w_j^\top x) = \operatorname{sgn}(w_{0,j}^\top x)$, i.e.,
$\sigma'(w_j^\top x) = \sigma'(w_{0,j}^\top x)$ almost surely.
Using the ReLU identity $f(x;W) - f_0(x;W) = \langle \nabla f(x;W) - \nabla f(x;W_0), W\rangle$
and two applications of Cauchy-Schwarz:
\[
  |f(x;W) - f_0(x;W)|
  \le \frac{1}{\sqrt{m}} \sum_{j \in S} \|w_j - w_{0,j}\|
  \le B \sqrt{\frac{|S|}{m}}
  \le \frac{2B^{4/3} + B\ln(1/\delta)^{1/4}}{m^{1/6}}.
\]
The second inequality is proved analogously by extending $S$ to include a third set
$S_3 := \{j : \|v_j - w_{0,j}\| \ge r\}$ and $S_4 = S_1 \cup S_2 \cup S_3$.
\end{proof}


\section{The neural tangent kernel}\label{sec:ntk:kernel}

The empirical NTK is defined as the Gram matrix of the gradient features.

\begin{lemma}[Frobenius inner product of gradients]\label{lem:ntk:frobInnerGradient}
  \uses{def:ntk:linearization}
  \lean{NTK.frobeniusInner_gradientMatrix_eq_empiricalNTKWithOuter}
The empirical NTK is exactly the Frobenius inner product of the network gradients:
\[
  k_m(x, x') = \langle \nabla_W f(x; W_0),\, \nabla_W f(x'; W_0) \rangle_{\mathrm{F}}.
\]
This factorization implies that $k_m$ is symmetric and positive semidefinite:
\end{lemma}

\begin{lemma}[Empirical NTK is positive semidefinite]\label{lem:ntk:empiricalNTKPSD}
  \lean{NTK.empiricalNTK_posSemidef}
For any finite set of points $x_1, \ldots, x_n \in \mathbb{R}^d$ and coefficients $\alpha_1, \ldots, \alpha_n \in \mathbb{R}$,
\[
  \sum_{i=1}^n \sum_{j=1}^n \alpha_i \alpha_j k_m(x_i, x_j) \ge 0.
\]
\end{lemma}

\begin{definition}[Empirical NTK]\label{def:ntk:empiricalNTK}
  \uses{def:ntk:shallowNetwork, def:ntk:gaussianInit, lem:ntk:frobInnerGradient}
  \lean{NTK.empiricalNTK}
Let $W_0 \in \mathbb{R}^{m \times d}$ and $a \in \{-1,+1\}^m$.
The \emph{empirical neural tangent kernel} at initialization $W_0$ is
\[
  k_m(x, x') := x^\top x' \cdot \frac{1}{m} \sum_{j=1}^{m}
    \sigma'(w_{0,j}^\top x)\, \sigma'(w_{0,j}^\top x'),
\]
where we have used $a_j^2 = 1$ and $\|x\|, \|x'\| \le 1$.
\end{definition}

\begin{definition}[Limiting NTK]\label{def:ntk:limitingNTK}
  \lean{NTK.limitingNTK}
The \emph{limiting neural tangent kernel} for activation $\sigma$ is
\[
  k(x, x') :=
  x^\top x' \cdot \mathbb{E}_{w \sim \mathcal{N}(0, I_d)}
    \!\bigl[\sigma'(w^\top x)\, \sigma'(w^\top x')\bigr].
\]
\end{definition}

\begin{lemma}[Almost sure convergence of the empirical NTK]\label{lem:ntk:ntkConverges}
  \uses{def:ntk:empiricalNTK, def:ntk:limitingNTK, def:ntk:gaussianInit}
  \lean{NTK.ntk_convergence}
For fixed $x, x' \in \mathbb{R}^d$, almost surely over
$W_0 \sim \mathcal{N}(0, I_d)^{\otimes m}$ (in $m$),
\[
  k_m(x, x') \xrightarrow{m \to \infty} k(x, x').
\]
\end{lemma}

\begin{proof}
  \uses{def:ntk:gaussianInit, def:ntk:empiricalNTK, def:ntk:limitingNTK}
The rows $(w_{0,j})_{j=1}^m$ are i.i.d.\ from $\mathcal{N}(0, I_d)$, so the summands
$\sigma'(w_{0,j}^\top x)\,\sigma'(w_{0,j}^\top x')$ are i.i.d. To apply the strong law of large numbers, we must ensure the summands are integrable.
\begin{lemma}[Integrability of the NTK summand]\label{lem:ntk:integrableSummand}
  \lean{NTK.integrable_ntkSummand}
If $\sigma'$ is measurable and bounded, the summand $w \mapsto \sigma'(w^\top x)\sigma'(w^\top x')$ is integrable with respect to the Gaussian row measure.
\end{lemma}
Because the ReLU derivative is bounded by 1, this integrability condition holds. 
The strong law of large numbers then gives the almost sure convergence of their empirical
average to the expectation $\mathbb{E}[\sigma'(w^\top x)\sigma'(w^\top x')]$.
\end{proof}

\begin{proposition}[ReLU NTK closed form]\label{prop:ntk:reluNTK}
  \uses{def:ntk:limitingNTK}
  \lean{NTK.reluNTK_closedForm}
For $\sigma = \sigma_{\mathrm{ReLU}}$ and $x, x' \in \mathbb{R}^d$ with $\|x\| = \|x'\| = 1$,
\[
  k(x, x') = x^\top x' \cdot \frac{\pi - \arccos(x^\top x')}{2\pi}.
\]
\end{proposition}

\begin{proof}
  \uses{def:ntk:limitingNTK}
Since $\sigma_{\mathrm{ReLU}}'(z) = \mathbf{1}[z \ge 0]$ a.e., we must compute
\[
  \mathbb{E}_{w \sim \mathcal{N}(0,I_d)}\!
    \bigl[\mathbf{1}[w^\top x \ge 0] \cdot \mathbf{1}[w^\top x' \ge 0]\bigr].
\]
By rotational invariance, $w^\top x$ and $w^\top x'$ depend only on the projection
of $w$ onto $\operatorname{span}(x, x')$; reducing to $\mathbb{R}^2$, the angle between
$x$ and $x'$ is $\theta := \arccos(x^\top x') \in [0,\pi]$.
The event $\{w^\top x \ge 0\} \cap \{w^\top x' \ge 0\}$ is a sector of angle $\pi - \theta$
(assuming $w$ uniform on the circle), so the probability is $(\pi - \theta)/(2\pi)$.
Multiplying by $x^\top x'$ gives the result.
\end{proof}


\section{Universal approximation via the NTK RKHS}\label{sec:ntk:universal}

\begin{definition}[NTK domain]\label{def:ntk:domain}
  \lean{NTK.ntkDomain}
The \emph{NTK domain} is the set
\[
  \mathcal{X} := \bigl\{x \in \mathbb{R}^d : \|x\| = 1,\; x_d = 1/\sqrt{2}\bigr\}.
\]
The last coordinate is fixed to $1/\sqrt{2}$ to serve as an implicit bias.
\end{definition}

To prove universality on this domain, we map the problem to a solid ball in $\mathbb{R}^{d-1}$.

\begin{lemma}[Reduced domain bijection]\label{lem:ntk:reducedDomain}
  \uses{def:ntk:domain, def:ntk:limitingNTK}
  \lean{NTK.ntkDomain_subset_sphere, NTK.isCompact_ntkDomain, NTK.reducedKernel_eq_reluNTK}
The NTK domain $\mathcal{X}$ is a compact subset of the sphere. By dropping the $d$-th coordinate (which is fixed to $1/\sqrt{2}$), $\mathcal{X}$ is in bijection with the reduced domain $U = \{u \in \mathbb{R}^{d-1} : \|u\|^2 \le 1/2\}$. The ReLU NTK on $\mathcal{X}$ equals the dot-product kernel $\tilde{k}(u, u') = \tilde{f}(u^\top u')$ on $U$, where
\[
  \tilde{f}(z) = \frac{z + 1/2}{2} - \frac{(z+1/2)\,\arccos(z+1/2)}{2\pi}.
\]
\end{lemma}

\begin{lemma}[Strictly positive Maclaurin coefficients]\label{lem:ntk:reducedKernelCoeff}
  \lean{NTK.reducedKernelCoeff_pos}
All Maclaurin series coefficients of $\tilde{f}(z)$ are strictly positive.
\end{lemma}

\begin{definition}[NTK RKHS predictor class]\label{def:ntk:rkhs}
  \uses{def:ntk:limitingNTK, def:ntk:domain}
  \lean{NTK.RKHSClass, NTK.zero_mem_RKHSClass, NTK.RKHSClass_smul, NTK.RKHSClass_add}
The \emph{NTK RKHS predictor class} (associated to the ReLU limiting NTK $k$) is the set
\[
  \mathcal{H} := \left\{
    x \mapsto \sum_{j=1}^{n} \alpha_j\, k(x, x_j)
    \;\Big|\;
    n \ge 0,\; \alpha_j \in \mathbb{R},\; x_j \in \mathcal{X}
  \right\}.
\]
\end{definition}

\begin{theorem}[Universal approximation via the NTK RKHS]\label{thm:ntk:universal}
  \uses{def:ntk:rkhs, def:ntk:domain, prop:ntk:reluNTK}
  \lean{NTK.isUniversal}
For $\sigma = \sigma_{\mathrm{ReLU}}$, the class $\mathcal{H}$ is a universal approximator
over $\mathcal{X}$: for every continuous $g : \mathcal{X} \to \mathbb{R}$ and every
$\epsilon > 0$, there exists $h \in \mathcal{H}$ such that
\[
  \sup_{x \in \mathcal{X}} |g(x) - h(x)| \le \epsilon.
\]
\end{theorem}

\begin{proof}
  \uses{def:ntk:rkhs, def:ntk:limitingNTK, prop:ntk:reluNTK, lem:ntk:reducedDomain, lem:ntk:reducedKernelCoeff}
Consider the bijection $\phi : \mathcal{X} \to U := \{u \in \mathbb{R}^{d-1} : \|u\|^2 \le 1/2\}$
that drops the $d$-th coordinate (Lemma~\ref{lem:ntk:reducedDomain}).
By the kernel universality criterion \cite{steinwart2008support} (Corollary~4.57),
a positive-definite kernel on a compact domain is universal if and only if its
associated power series (in the dot-product case) has all-positive coefficients.
Because all Maclaurin coefficients of $\tilde{f}$ are strictly positive (Lemma~\ref{lem:ntk:reducedKernelCoeff}), $\tilde{k}$ is a universal kernel over $U$, and thus $\mathcal{H}$ is a universal
approximator over $\mathcal{X}$.
\end{proof}

Finally, by combining the universal approximation of the RKHS class with the ReLU linearization bound, we show that wide neural networks themselves can approximate any continuous function on the NTK domain near initialization.

\begin{theorem}[Network approximation near initialization]\label{thm:ntk:networkApprox}
  \uses{def:ntk:rkhs, lem:ntk:reluLinearization}
  \lean{NTK.rkhs_approx_by_network}
For every RKHS function $h \in \mathcal{H}$, target accuracy $\epsilon > 0$, and failure probability $\delta \in (0, 1)$, there exists a large enough width $m$, a radius $B \ge 0$, and an outer layer $a \in \mathbb{R}^m$ such that with probability at least $1-\delta$ over Gaussian initialization $W_0$, there exists an inner weight matrix $W$ with $\|W-W_0\|_{\mathrm{F}} \le B$ satisfying:
\[
  \sup_{x \in \mathcal{X}} |f(x; W) - h(x)| \le \epsilon.
\]
\end{theorem}
